% !TeX root = ../../Book3.tex
% 纲要标题：## 附录 B　计算交换代数专题
% 纲要位置：工作记录/计划纲要.md:2857

\chapter{计算交换代数专题}
\label{b3:app:B}

\begin{chapterabstract}
本附录在明确输入表示和计算证书的前提下，讨论零维理想的 FGLM 换序、模素数计算、参数分支及单项式计数。模 Gröbner 基与合冲构造用于保存生成元变换，并核验自由消解的正合性；各项算法的执行都要求能够精确运算与判零。
\end{chapterabstract}

\begin{learninggoals}
\begin{itemize}
\item 明确系数环、工作环、项序、精度和展示矩阵在计算中的作用。
\item 用商代数的乘法矩阵执行 FGLM 换序，证明算法终止并核验结果。
\item 区分一般参数与退化分支，用有限等式验证重构与分解的候选答案。
\item 追踪理想和模的生成元变换，分清矩阵复形与已经证明正合的消解。
\end{itemize}
\end{learninggoals}

在 \(k[x,y]\) 中，理想 \(I=(x^2-1,y-x)\) 给出的商由 \(1,x\) 张成，且 \(x^2=1\)、\(y=x\)。因此乘以 \(x\) 在这两个基向量上的矩阵可以直接写出。若改变变量次序或消去目标，这个有限维商不变，但适合计算的新 Gröbner 基可能完全不同；由乘法矩阵恢复新关系，是换序方法的起点。

\ProseParagraphBreak

计算结束后，除了一个“答案”，还需要能检查它的资料：理想成员应附生成元表示，不属于理想要有适用的正规式，模素数计算要有恢复到原系数域的依据。零维性使商代数具有有限线性维数，因而允许在有限矩阵中转移问题；参数或局部情形则可能改变这个前提，须另外分支处理。

\ProseParagraphBreak

本附录主要使用\chapref{b3:ch:08}的全局 Gröbner 基与零维商代数；分解部分还用到\chapref{b3:ch:04}的准素分解，模计算部分使用\chapref{b3:ch:21}。局部标准基的具体方法另在\appref{b3:app:F}讲述。

FGLM 在这里给出完整步骤、终止性和正确性证明；局部成员等式与正规化条件也给出完整认证论证。F4 的矩阵化、签名重写和正则链主要用来解释运算机制与输出对象。F5、Gröbner walk、综合 Gröbner 系统的构造及正则链分解还需要各自的候选管理、终止性与完整性理论，本文只说明所用思想和查阅依据。

\ProseParagraphBreak
Cox、Little 与 O'Shea 的《Using Algebraic Geometry》从消去方程与商环结构的配合引入计算问题，见\cite[Chapter 2, pp. 24--25]{CoxLittleOShea1998UsingAG}。FGLM 把换序转为商代数中的线性相关计算，F4 型方法把多项式约化组织为稀疏矩阵消元；参数分支与模素数计算则分别处理特殊化和系数增长带来的困难。这些计算虽然形式不同，最终都要在指定的环中核验结果。

设系数域 \(k\) 可进行精确的四则运算与判零，\(S=k[x_1,\ldots,x_n]\)。\cref{b3:tab:B-computation-certificates}比较几类任务的输入条件与核验证据；其中 \(F\) 表示输入生成组，模计算时表示生成矩阵。

\begin{table}[htbp]
\centering
\caption{计算任务、输入模型与核验证据}
\label{b3:tab:B-computation-certificates}
\begin{tabularx}{\linewidth}{@{}>{\raggedright\arraybackslash}p{0.12\linewidth}>{\raggedright\arraybackslash}p{0.25\linewidth}>{\raggedright\arraybackslash}p{0.30\linewidth}>{\raggedright\arraybackslash}X@{}}
\toprule
任务 & 输入与计算 & 核验证据 & 候选本身尚未证明的性质 \\
\midrule
理想成员判定 & 生成组 \(F\)、\(f\) 与全局项序；使用 Gröbner 基及多元除法 & \(f=Fa\) 证明成员；已验证基的非零正规余式证明非成员 & 一份像基的多项式列表尚未证明生成原理想 \\
换序 & 已知 Gröbner 基、新项序与有限维商；使用 FGLM & 商代数坐标、候选处理记录及各步线性相关等式 & 新关系列表尚未证明没有遗漏标准单项式 \\
有理系数计算 & 有理多项式生成组；使用模素数计算及重构 & 在 \(\mathbb Q\) 上精确核对两向表示与临界对标准表示 & 多个素数下形状一致尚未证明有理候选正确 \\
参数分支 & 参数方程；按综合 Gröbner 系统分支 & 覆盖条件及各分支上可特殊化的基与表示 & 一般参数处的答案尚未覆盖特殊纤维 \\
合冲与消解 & 多项式矩阵及基次数；使用模 Gröbner 基 & 全部核的生成矩阵、生成元变换及像等于核的证书 & \(d_id_{i+1}=0\) 只证明复形，尚未证明正合 \\
\bottomrule
\end{tabularx}
\end{table}

表中的计算方法负责产生候选，核验证据负责把候选落实为所声称的数学对象。输入环、项序或精度改变时，相应证据也须在新的工作环中重新解释。

% !TeX root = ../../Book3.tex
% 纲要标题：### B.1 系数、表示与复杂度

\section{系数、表示与复杂度}
\label{b3:sec:B01}

计算理想或商环，首先需要把对象写成有限数据，并说明怎样检验输出是否代表原来的对象。本节以多项式、矩阵和分式为例，区分可直接执行的运算与仍须给出证明的结构判断。

% 纲要细目（与计划纲要同步）：
% * 多项式、局部化分式、幂级数截断和有限展示模的不同输入模型
% * 域上的 Gröbner 基与整环、主理想整环上的强 Gröbner 基；首系数的额外作用
% * 有效系数环所需的可判定性和运算条件；不声称任意 Noether 环都支持有效算法
% * 算术复杂度、次数增长、系数膨胀和最坏情形与常见实例的区别
%
% #### B.1.1 多项式、分式、幂级数截断与有限展示模的表示
% #### B.1.2 Gröbner 基、强 Gröbner 基与首系数
% #### B.1.3 有效系数环的可判定性与运算条件
% #### B.1.4 算术复杂度、次数增长与系数膨胀


\subsection{多项式、分式、幂级数截断与有限展示模的表示}
\label{b3:subsec:B01:01}

理想写成有限生成形式，并不等于已经能判定成员关系；有限维商环写出一个向量空间基，也还没有给出乘法。计算之前应先说明，输入包含哪些有限数据，允许执行哪些运算，答案又以什么形式返回。

\ProseParagraphBreak
以下以可进行精确运算的域 \(k\) 为系数域。多项式用非零系数与指数向量的有限表表示；稀疏表示保存非零项，稠密表示则保存给定次数范围内的全部系数。局部化 \(U^{-1}k[x]\) 的元素用一对多项式 \((f,u)\) 表示，其中 \(u\in U\)。这对数据还应说明分母属于哪个乘性集：在原点局部环中，条件是 \(u(0)\ne0\)；在 \(k[x]_h\) 中，分母可取 \(h\) 的幂。这两种分式不能仅凭同样的斜线记号混用。

\ProseParagraphBreak
任意形式幂级数通常没有有限的输入表。有限计算可以接收一个多项式截断、一个逐次提供系数的规则，或者满足某个方程并附带分支条件的特殊级数。前两者所能证明的结论不同：知道
\[
f\equiv f_N\pmod{(x_1,\ldots,x_n)^N}
\]
只确定次数小于 \(N\) 的系数，不能据此断言 \(f=0\)。例如零级数和 \(x_1^N\) 在这份数据中不可区分。

\begin{definition}\label{b3:def:B-presentation-data}
设 \(R\) 为具有指定有限数据表示的交换环，\(p,q\) 为非负整数。给定矩阵 \(A\in R^{q\times p}\)，本附录以
\[
R^p\xrightarrow{A}R^q\longrightarrow M\longrightarrow0
\]
作为模 \(M=\operatorname{coker}A\) 的展示数据；矩阵的列是关系，\(R^q\) 的标准基给出指定生成元。若 \(M\) 有分次，还须给出源、靶各基向量的次数，使每个矩阵项的次数与所表示的映射相符。
\end{definition}

给出 \(\operatorname{coker}A\) 无须断言左映射单射。两个展示描述同构模时，计算仍须记录生成元如何对应，否则从一个展示求出的关系不能直接当作另一个展示的关系。这一点在\secref{b3:sec:B04}中用矩阵等式具体说明。

\subsection{Gröbner 基、强 Gröbner 基与首系数}
\label{b3:subsec:B01:02}

域上的除法能把任意非零首系数变成 \(1\)。在整数上，这一步可能改变理想。设 \(R\) 为整环，并在 \(R[x_1,\ldots,x_n]\) 的单项式上固定全局项序。首单项式 \(\operatorname{LM}(f)\) 只记录幂；首项 \(\operatorname{LT}(f)\) 则包含系数。若
\[
\operatorname{LT}(g)=a x^\alpha,\qquad
\operatorname{LT}(f)=b x^\beta,
\]
用 \(g\) 一步消去 \(f\) 的首项，既要求 \(x^\alpha\mid x^\beta\)，又要求 \(a\mid b\)。

\begin{definition}\label{b3:def:B-strong-groebner}
设 \(R\) 为整环，\(I\subset R[x_1,\ldots,x_n]\) 为理想，项序为全局项序。若有限集 \(G\subset I\setminus\{0\}\) 满足：每个非零 \(f\in I\) 的首项都被某个 \(g\in G\) 的首项整除，则称 \(G\) 为 \(I\) 的\term{强 Gröbner 基}[strong Gröbner basis]。这里的整除同时包括系数整除和单项式整除。
\end{definition}

另一种较弱的条件是，所有 \(\operatorname{LT}(g)\) 生成 \(I\) 的首项理想。这两种条件在域上相同，在一般系数环上不同。例如在 \(\mathbb Z[x]\) 中，\(\{2x,3x\}\) 的首项生成 \((x)\)，但 \(x=3x-2x\) 的首项既不被 \(2x\) 整除，也不被 \(3x\) 整除。因此它不是 \((x)\) 的强 Gröbner 基，加入 \(x\) 后才得到单项首项整除的性质。

\ProseParagraphBreak
在有效主理想整环上，首系数的最大公因子及 Bézout 表示提供了额外的构造：若两个首项有系数 \(a,b\)，就可用 \(ua+vb=\gcd(a,b)\) 产生系数更小的首项。仅计算域上的 \(S\)-多项式会漏掉这一要求。本附录不建立整环上强 Gröbner 基的完整算法；后文的 FGLM、矩阵消元和签名说明均以域为系数环。

\begin{example}\label{b3:exm:B-integer-coefficients}
在 \(\mathbb Z[x]\) 中取 \(I=(2x,x^2)\)。集合 \(\{2x,x^2\}\) 是强 Gröbner 基：\(I\) 中次数为 \(1\) 的多项式，其 \(x\) 系数为偶数；次数至少为 \(2\) 的非零多项式，其首项被 \(x^2\) 整除。但 \(x\notin I\)。扩张到 \(\mathbb Q[x]\) 后却有 \(I\mathbb Q[x]=(x)\)。因此，先在分式域上求基再清除分母，不能自动恢复原整数理想的成员关系。
\end{example}

\subsection{有效系数环的可判定性与运算条件}
\label{b3:subsec:B01:03}

“系数环是 Noether 环”保证某些生成过程在抽象意义上有限，却没有告诉我们如何比较两个输入系数。有效算法首先需要有限可读的元素表示、相等判定、加法与乘法；在域上还需要非零元素的求逆。在一般系数环上，约化可能进一步需要判定一个系数是否属于给定的有限生成系数理想，并在属于时求出线性组合。

\ProseParagraphBreak
例如有理数可用互素整数对表示；整数的精确四则运算和最大公因子计算给出相等判定。有限域须指定其特征以及扩域时所用的不可约多项式。代数数域可以给出 \(\mathbb Q[T]/(p)\) 及 \(p\) 的不可约性依据，再把系数约化到次数小于 \(\deg p\)。若所关心的是实根中的某个嵌入，还需隔离区间等附加数据，但抽象域上的理想运算并不依赖选哪一个复嵌入。

\ProseParagraphBreak
浮点数的近似零判定不具备同样性质。一个很小的非零主元可以决定矩阵秩，也可以决定某个多项式的首项。精确 Gröbner 基证书需要确定这些系数是否为零；数值计算则须另有误差控制和认证方法。两类计算可以配合，但不能把“数值余项很小”作为理想成员的精确证明。

\begin{remark}
若输入系数是任意实数，而所谓表示仅允许不断给出近似值，则通常无法在有限时间内判定它是否恰为零。因此，“任意域上的定理”与“任意域上的可执行程序”是不同层次的陈述。下文说算法可执行时，始终包含了前述精确系数运算假设。
\end{remark}

\subsection{算术复杂度、次数增长与系数膨胀}
\label{b3:subsec:B01:04}

计算代价至少有三种来源：系数运算的次数、中间多项式出现的单项式数，以及每个系数所需的存储位数。以域运算计数，可以比较两种线性代数方法；若系数在 \(\mathbb Q\) 中，还必须考虑分子、分母的增长。同样次数和同样项数的多项式，系数为一位整数与为数千位整数时，实际计算量并不相近。

\begin{example}\label{b3:exm:B-bit-growth}
在 \(\mathbb Q[x_1,\ldots,x_r]\) 中考虑
\[
x_1-2,\quad x_2-x_1^2,\quad\ldots,\quad x_r-x_{r-1}^2.
\]
每个输入的次数至多为 \(2\)，系数很小，但唯一解为
\[
x_i=2^{\,2^{i-1}}.
\]
若把消元结果写为 \(x_r-2^{2^{r-1}}\)，常数项的二进制位数是 \(2^{r-1}+1\)。输出本身已要求指数增长的位数，因而不能仅凭输入次数小就期待小规模的显式答案。
\end{example}

次数增长也会扩大线性代数问题。\(n\) 个变量中次数恰为 \(d\) 的单项式有 \(\binom{n+d-1}{d}\) 个，次数至多为 \(d\) 的有 \(\binom{n+d}{d}\) 个。用一张系数矩阵同时处理许多约化时，这些数给出可能的列数。稀疏存储能避开最初不存在的系数，但行消元可能产生新的非零项。

\ProseParagraphBreak
因此，算法说明应把可证明的终止性、给定模型下的运算界和实际实例的表现分别报告。F4 型方法利用批量约化，FGLM 利用商代数已经有限维，模方法则控制有理系数的膨胀；它们改善的是不同部分，没有一种选择对所有输入都最好。

% !TeX root = ../../Book3.tex
% 纲要标题：### B.2 Gröbner 基的现代计算策略与 FGLM 换序
% 纲要位置：工作记录/计划纲要.md:2871

\section{Gröbner 基的现代计算策略与 FGLM 换序}
\label{b3:sec:B02}

逐项求余会重复消去许多相同单项式，而系数增长也可能使精确运算变慢。把约化组织成矩阵、转到有限域计算再重构，都需要保留能够回到原系数环核验的证据。

% 纲要内容（仅作写作计划，尚非教材正文）：
%
% * F4 的符号预处理与稀疏消元；签名是指定模序的首单项式，重写机制不等于完整 F5 算法
% * FGLM 的接受项边界探索、候选数界及正确性四段证明；错误扩展反例在正文证明；Gröbner walk 说明思想与依据
% * 模素数形状分组、CRT、有理重构及原域精确验证；固定证书的有限坏素数集是事后结论
% * 两向表示负责理想相等，临界对标准表示负责首项条件；完整有限证书算例与模、消解证书分工
% #### B.2.1 F4、F5 型算法、签名准则与稀疏线性代数
% #### B.2.2 FGLM 零维换序与 Gröbner walk
% #### B.2.3 模方法、坏素数、重构与原域验证
% #### B.2.4 算法实现与结果验证
%
% ---
%


\subsection{F4、F5 型算法、签名准则与稀疏线性代数}
\label{b3:subsec:B02:01}

Buchberger 算法逐对形成 \(S\)-多项式并求余，其中许多约化会反复使用同一个单项式倍数。F4 型方法\cite{Faugere1999F4}把共享约化行的一批首项消去组织成同一个稀疏线性代数问题：先选取临界对，收集所需多项式倍数；以出现的单项式为列、这些倍数的系数为行；再作行消元，提取产生新首项的行。数学上仍在消去首项，改变的是这些运算的组织方式。每个行变换都是系数域上的可逆线性变换，因此保留这批多项式张成的向量空间，也保留它们属于原理想的表示。

\ProseParagraphBreak
这里还有一个不可省略的步骤：收集用于消去首项的约化行。例如一行出现了可被现有首项整除的单项式，就应加入相应的基元素倍数，再检查新引入的项。这称为符号预处理。矩阵必须同时容纳候选 \(S\)-多项式、所需的约化行以及这些约化行带来的新单项式列；单独求若干多项式行的行最简形，只完成了一个向量空间的计算。临界对的选择、符号预处理以及尚未处理的对仍需统一管理，才能得到整个理想的 Gröbner 基。这里说明矩阵约化的机制，完整 F4 算法的临界对管理和正确性证明见所引文献。

\ProseParagraphBreak
例如在 \(\mathbb Q[x,y]\) 中取次数字典序，\(x>y\)，当前多项式为 \(g_1=x^2-y\)、\(g_2=xy-1\)、\(g_3=y^2-y\)。选取前两者的临界对，得到
\[
S=yg_1-xg_2=x-y^2.
\]
若只放入这一行，单项式列是 \(y^2,x\)，却不能完成对首项 \(y^2\) 的约化。符号预处理发现它被 \(\operatorname{LM}(g_3)=y^2\) 整除，补入约化行 \(g_3\)，并把新出现的 \(y\) 加为一列。按 \(y^2,x,y\) 排列的矩阵及一次消元为
\[
\begin{array}{c|rrr}
&y^2&x&y\\\hline
S&-1&1&0\\
g_3&1&0&-1
\end{array}
\quad\longrightarrow\quad
\begin{pmatrix}1&0&-1\\0&1&-1\end{pmatrix}.
\]
右侧先保留 \(g_3\) 的行，再把两行相加，得到新行 \(x-y=S+g_3\)。每一行都来自原理想 \((g_1,g_2,g_3)\)，且新首项 \(x\) 不被原来的 \(x^2,xy,y^2\) 整除。这显示预处理为何要补约化行和新列；这一批的消元仍不代替对其余临界对的检查。

\begin{definition}\label{b3:def:B-signature}
设 \(k\) 为域，\(S=k[x_1,\ldots,x_n]\)，取 \(f_1,\ldots,f_s\in S\)，以 \(e_1,\ldots,e_s\) 记 \(S^s\) 的标准基，并在其上指定全局模项序。对于带有表示
\[
(u,f),\qquad u=(u_1,\ldots,u_s),\quad f=\sum_i u_if_i,\quad u\ne0
\]
的数据，\(\operatorname{LM}(u)\) 是一个带基向量指标的模单项式 \(x^\alpha e_i\)，称为这份表示的\term{签名}[signature]。签名依赖指定的模项序；同一个多项式的不同表示可以有不同签名。
\end{definition}

例如 \(f_1=x,f_2=y\) 时，\(y e_1-x e_2\) 表示零多项式，却携带一个非零模首项。若模序先比较基分量，规定 \(e_1>e_2\) 时其签名为 \(y e_1\)，规定 \(e_2>e_1\) 时其签名为 \(x e_2\)。所以签名记录的是生成元表示的模首单项式，与矩阵秩、数值指纹或软件散列值不同。下面的命题说明如何利用一条已知关系降低表示的签名，是签名方法的基本代数机制。

\begin{proposition}\label{b3:prop:B-signature-rewrite}
设 \(k\) 为域，\(S=k[x_1,\ldots,x_n]\)，取 \(f_1,\ldots,f_s\in S\)，令 \(F:S^s\to S\) 为 \((u_i)\mapsto\sum_i u_if_i\)，并固定 \(S^s\) 上的全局模项序。若非零 \(z\in\ker F\) 的首单项式整除非零向量 \(u\in S^s\) 的首单项式，则 \(Fu\) 有一个向量表示 \(u'\)，满足 \(u'=0\) 或 \(\operatorname{LM}(u')<\operatorname{LM}(u)\)。
\end{proposition}
\begin{proof}
两个首单项式的整除意味着它们在同一基分量中，故可取一个带系数单项式 \(c x^\alpha\)，使 \(u\) 与 \(c x^\alpha z\) 的首项相同。令 \(u'=u-cx^\alpha z\)，首项消去后得到所需严格不等式，而 \(Fu'=Fu\)，因为 \(Fz=0\)。
\end{proof}

F5 算法\cite{Faugere2002F5}及其他签名算法把关系信息用于筛选候选表示，并限制约化步骤，使签名信息能够继续追踪。上面的命题只给出了降低签名的一种操作；若据此省去某份候选，算法还须保证较低签名的代表会被处理。完整的 F5 算法还需要规定临界对、允许的约化、重写规则及处理次序，并分别证明筛选后的完备性、终止性与正确性。这些内容不由该命题推出，本节只介绍其基本机制。最终得到的候选 Gröbner 基可用\subsecref{b3:subsec:B02:04}中的普通证书核验：检查表示矩阵和临界对的约化记录，即可检验其生成的理想与首项条件。

\subsection{FGLM 零维换序与 Gröbner walk}
\label{b3:subsec:B02:02}

\subsecref{b3:subsec:0805:03}已把换序解释为商代数中的线性相关判定。FGLM 算法将这一思路用于有限维商代数\cite{FaugereGianniLazardMora1993FGLM}，它同样适用于非根理想和非代数闭的系数域。

\ProseParagraphBreak
算法维护已经接受的单项式表 \(B\)，只从每个 \(b\in B\) 产生相邻倍数 \(x_1b,\ldots,x_nb\)。这样便沿着已接受单项式集合的边界向外探索：候选的商代数坐标若与已保存坐标线性无关，就加入 \(B\) 并继续产生后继；若相关，就得到以该候选为首项的关系，并停止从它产生后继。队列只在这些候选中取目标序的最小元，无须枚举目标序下的全部单项式。

\begin{algorithm}[htbp]
\caption{FGLM 换序}\label{b3:alg:B-fglm}
\begin{algorithmsteps}
{精确可计算域 \(k\)，真理想 \(I\subset S=k[x_1,\ldots,x_n]\) 的一个全局 Gröbner 基，满足 \(D=\dim_k(S/I)<\infty\)，以及目标全局项序 \(<\)。}
{理想 \(I\) 关于 \(<\) 的 Gröbner 基 \(G\)，以及新标准单项式表 \(B\)。}
\State 在旧标准单项式基中计算乘法矩阵 \(T_i\)，记 \(v(m)\in k^D\) 为单项式 \(m\) 的坐标。
\State 置有序候选集合 \(Q=\{1\}\)，已处理集合 \(P=\varnothing\)，并置 \(B=G=\varnothing\)。
\While{\(Q\ne\varnothing\)}
\State 取出 \(Q\) 中最小的单项式 \(m\)，将 \(m\) 记入 \(P\)。
\If{\(m\) 被某个 \(g\in G\) 的首单项式整除}
\State 丢弃 \(m\)，继续取下一个候选。
\Else
\State 计算 \(v(m)\)，判定它与已保存的 \(v(b)\ (b\in B)\) 是否线性相关。
\If{\(v(m)=\sum_{b\in B}c_b\cdot v(b)\)}
\State 把 \(m-\sum_{b\in B}c_bb\) 加入 \(G\)。
\Else
\State 把 \(m\) 加入 \(B\)，并把尚未属于 \(Q\cup P\) 的 \(x_1m,\ldots,x_nm\) 加入 \(Q\)。
\EndIf
\EndIf
\EndWhile
\State 返回 \(G,B\)。若需要约化 Gröbner 基，最后作首一化与交互约化。
\end{algorithmsteps}
\end{algorithm}

只从加入 \(B\) 的单项式产生后继，是下述候选数估计所需的算法条件。坐标可由 \(v(x_i b)=T_i\cdot v(b)\) 计算，也可用旧 Gröbner 基求余；保存坐标可以避免重复运算。若 \(I=S\)，旧 Gröbner 基已检出 \(1\)，直接返回 \(\{1\}\) 和空基即可，不进入上述真理想版本。

\begin{theorem}\label{b3:thm:B-fglm}\label{b3:thm:B-fglm-candidate-bound}
设 \(k\) 是支持精确四则运算和零判定的域，\(I\subset S=k[x_1,\ldots,x_n]\) 为真理想，且 \(D=\dim_k(S/I)<\infty\)。给定 \(I\) 关于一个可有效比较的全局项序的 Gröbner 基，并指定另一可有效比较的全局项序 \(<\)。在\cref{b3:alg:B-fglm}中，后继只由加入 \(B\) 的单项式产生，拒绝的候选不再扩展。因此算法处理的候选总数至多为 \(1+nD\)，并在有限步后返回 \(S/I\) 的目标标准单项式基 \(B\) 及 \(I\) 关于 \(<\) 的 Gröbner 基 \(G\)。
若改为从每个处理过的单项式产生后继，包括被首项整除判据拒绝的候选，则终止性一般失效：在 \(k[x]/(x)\) 中可以产生永不清空的 \(x,x^2,x^3,\ldots\) 队列。
\end{theorem}
\begin{proof}
每次加入 \(B\) 时，坐标向量与原有向量线性无关，故这种加入至多发生 \(D\) 次。除初始候选 \(1\) 外，每个候选都由某次加入 \(B\) 产生，而每次加入至多产生 \(n\) 个候选，因而总共至多有 \(1+nD\) 个候选。每个候选只处理一次，所以算法终止。这里不能把产生后继的步骤移到每次候选处理之后：在 \(k[x]/(x)\) 中，接受 \(1\) 后，\(x\) 的零坐标给出关系 \(x\)；若仍从被拒绝的 \(x^r\) 产生 \(x^{r+1}\)，即使所有后来的候选都立即被首项整除判据丢弃，队列仍会永远非空。

全局项序满足 \(m<x_i m\)。由队列取最小元和只从已处理单项式产生更大单项式可知，处理次序严格递增。在线性相关时，关系
\[
g=m-\sum_{b\in B}c_bb
\]
中所有 \(b\in B\) 都是此前已经处理并接受的单项式，故 \(b<m\)。右端的较小单项式无法消去 \(m\)，因此 \(\operatorname{LM}(g)=m\)；它在旧标准单项式基中的坐标为零，故 \(g\in I\)。最终 \(B\) 的元素都不被 \(G\) 的首单项式整除：在它加入时不被已有首项整除，以后产生的首项又比它大，不可能整除它。

反过来，凡不被这些首单项式整除的单项式都在 \(B\) 中。否则取这样的最小单项式 \(m\)。\(m\ne1\)，因为真理想情形下 \(v(1)\ne0\)，于是 \(1\) 被接受。写 \(m=x_i b\)，其中 \(b\) 是一个真因子。\(b\) 也不被任何首单项式整除，且 \(b<m\)，故由最小性 \(b\in B\)。于是 \(m\) 已进入队列。处理它时，既不能按首项整除丢弃，也不能因相关而产生以 \(m\) 为首项的关系，只能加入 \(B\)，矛盾。

令 \(J=(G)\subset I\)。全局多项式除法把每个 \(f\in S\) 写成 \(G\) 的组合加上由 \(B\) 张成的余项。若 \(f\in I\)，这个余项在 \(S/I\) 中为零；而 \(B\) 的坐标线性无关，故余项为零，得到 \(I=J\)。若非零 \(f\in I\) 的首单项式不被任何 \(\operatorname{LM}(g)\) 整除，除法会把它保留为余项的最高项，后续较小项无法消去它，与余项为零矛盾。因此 \(G\) 是 Gröbner 基。由余项的张成性及坐标独立性，\(B\) 是商空间基，也是目标标准单项式集合。
\end{proof}

\begin{example}\label{b3:exm:B-fglm-complete}
取 \(I=(x-y^2,y^3-2)\subset\mathbb Q[x,y]\)，旧序为 \(x>y\) 的字典序，旧 Gröbner 基为 \(\{x-y^2,y^3-2\}\)，旧标准单项式基为 \((1,y,y^2)\)。在后一个基给出的坐标中
\[
T_y=\begin{pmatrix}0&0&2\\1&0&0\\0&1&0\end{pmatrix},
\qquad T_x=T_y^2.
\]
改用 \(y>x\) 的字典序。本次实际进入队列的候选按该序排列为
\[
1<x<x^2<x^3<y<xy<x^2y.
\]
接受 \(1\) 后，候选为 \(x,y\)；先取 \(x\)，其坐标为 \((0,0,1)^{\mathsf T}\)，接受。接着取 \(x^2\)，其坐标为 \((0,2,0)^{\mathsf T}\)，仍接受。再取 \(x^3\)，其坐标为 \((4,0,0)^{\mathsf T}\)，所以产生 \(x^3-4\)，不再从 \(x^3\) 产生后继。虽然 \(x^r<y\) 对所有 \(r\ge1\) 都成立，\(x^4,x^5,\ldots\) 不会进入这个队列。

\ProseParagraphBreak
随后取 \(y\)，因 \(v(y)=\dfrac12v(x^2)\)，产生 \(y-\dfrac12x^2\)。已在队列中的 \(xy,x^2y\) 被首项 \(y\) 整除，全部丢弃。最后
\[
B=(1,x,x^2),\qquad
G=\{x^3-4,\ y-\dfrac12x^2\}.
\]
两个关系都可在旧商环直接核验。它们的首项 \(x^3,y\) 给出恰好三个标准单项式，与算法保存的独立坐标一致。
\end{example}

有限维是上述队列计数的根据。对正维理想，标准单项式无限，不能照搬终止证明。作为另一种换序机制，Gröbner walk\cite{CollartKalkbrenerMall1997GroebnerWalk}用权向量及必要的细化次序描述项序：在一个首项不变的区域内移动权向量，跨过区域边界时计算相应权初始理想，再把关系提升回原理想。这一段只说明换序的几何思路；完整算法还需定义权初始理想与 Gröbner 扇，给出跨越边界的规则并证明提升步骤的正确性，详见所引文献。FGLM 的完整算法与证明则以上面的有限维商代数计算为依据。

\subsection{模方法、坏素数、重构与原域验证}
\label{b3:subsec:B02:03}

这里的模素数方法（modular methods）指先对系数取模再计算，中文的“模”是模素数的取模运算，与后文的模及合冲计算有别。在 \(\mathbb Q[x]\) 中计算时，可以先清除输入分母，在有限域上试算，再尝试恢复有理系数。这样做减少了中间系数膨胀，但必须区分三件事：某个素数下的计算正确，若干素数下的答案一致，以及重构答案在 \(\mathbb Q\) 上正确。

\ProseParagraphBreak
坏素数可能使首系数变零、分母不可逆，或使某个秩下降。例如输入 \(2x+y,y^2-1\) 在特征不为 \(2\) 时给出 \(x=-\dfrac y2\)，而模 \(2\) 后成为 \(y,y^2-1\)，生成整个环。这个素数下的正确答案不能描述原有的零维商代数。

\begin{proposition}[证书给出的有限坏素数集]\label{b3:prop:B-good-primes-certificate}
设 \(F\) 为 \(\mathbb Q[x_1,\ldots,x_n]\) 的一个有限多项式列，固定全局项序，\(G\) 是同一理想的首一 Gröbner 基。固定 \(G\) 由 \(F\) 的表示、\(F\) 由 \(G\) 的表示，以及所有 \(S\)-多项式约化为零的有限除法记录。存在非零整数 \(N\)，使这些记录可在 \(\mathbb Z[1/N]\) 中书写，且对每个 \(p\nmid N\)，模 \(p\) 后的 \(G\) 仍为模 \(p\) 后 \(F\) 的 Gröbner 基，并保留首单项式。这里的 \(N\) 从已经固定的有理基元素、两向表示和约化记录的系数中构造，给出的是这份证书的事后保证。
\end{proposition}
\begin{proof}
取 \(N\) 为全部输入、基元素及有限记录所出现系数的分母的公倍数，并包含使各首项标准化所用的非零整数因子。模 \(p\nmid N\) 时，所有等式有意义。首一性使基的首项不消失；两向表示证明两列生成同一理想。每份除法记录给出的 \(S\)-多项式表示，其非零乘积的首单项式严格小于对应的首单项式最小公倍数。取模只会删去项，基元素的首单项式及其最小公倍数不变，因而这些严格上界仍成立。由\cref{b3:thm:buchberger-criterion}，模 \(p\) 后的基仍为 Gröbner 基，标准单项式集合也随之保持。
\end{proof}

该命题所给的可能坏素数集是 \(N\) 的素因子集；它由正确的有理证书确定，其中某些素数也可能保持同样的首项。在有理答案尚未知时，这个命题不能事先找出全部好素数。实际计算时，先在多个有限域上计算，按首单项式形状分组，再利用中国剩余定理合并同形答案，作有理重构，最后回到原域精确验证。前面的计算和重构只帮助产生有理候选，最后的验证负责确认它的正确性。若 \(M\ge2\)，模 \(M\) 的剩余类 \(c\) 要在给定整数界 \(A\ge0\)、\(B\ge1\) 内恢复为 \(\dfrac ab\)，可要求
\[
a\equiv cb\pmod M,\quad |a|\le A,\quad 0<b\le B,\quad
\gcd(a,b)=1,\quad \gcd(b,M)=1.
\]
最后一个条件使分母模 \(M\) 可逆。当 \(2AB<M\) 时，范围内答案至多一个：若 \(a/b\) 与 \(a'/b'\) 都满足这些条件，则 \(ab'-a'b\) 被 \(M\) 整除，且其绝对值至多为 \(2AB<M\)，所以交叉差为零。分母为正且分子分母互素，又使这两个分数的表示相同。合理的大小界未必预先可知，重构失败时应继续加入素数。

\ProseParagraphBreak
最终的有理候选仍须回到原系数域核验两向生成关系与 Gröbner 条件。许多素数下结果一致可以帮助选择候选，却不能替代这一步；验证成功后，结论就不依赖猜测哪些素数是好的。

\subsection{算法实现与结果验证}
\label{b3:subsec:B02:04}

本小节检查一个理想生成组及其 Gröbner 条件；模计算中核是否全部求出、复形是否正合以及消解是否极小，见\subsecref{b3:subsec:B04:01}。设 \(k\) 为域，\(S=k[x_1,\ldots,x_n]\)，固定全局项序，输入行向量为 \(F=(f_1,\ldots,f_s)\)。删去候选中的零元素后，记候选行向量为 \(G=(g_1,\ldots,g_t)\)。一份可单独检查的证书应包含工作域、变量和项序，以及下列有限数据：
\begin{enumerate}
\item 一个矩阵 \(U\in S^{s\times t}\)，使 \(G=FU\)。这给出 \((G)\subseteq(F)\)，证明每个候选仍在输入理想中。
\item 一个矩阵 \(V\in S^{t\times s}\)，使 \(F=GV\)，即每个 \(f_i=\sum_jv_{ji}g_j\)。这给出 \((F)\subseteq(G)\)，证明没有丢掉原理想。
\item 对每个 \(1\le p<q\le t\)，一个表示 \(S(g_p,g_q)=\sum_jh_jg_j\)，其中 \(h_j\in S\)，且每个非零乘积满足
\[
\operatorname{LM}(h_jg_j)<
\operatorname{lcm}\bigl(\operatorname{LM}(g_p),\operatorname{LM}(g_q)\bigr).
\]
\end{enumerate}
前两项给出 \((F)=(G)\)。第三项是具有严格首项上界的\term{标准表示}[standard representation]；按 \(G\) 除法所得的零余项记录会给出这种表示，而任意等式 \(S=\sum_jh_jg_j\) 未必具有该上界。由\cref{b3:thm:buchberger-criterion}的有界表示条件，第三项证明 \(G\) 确为该项序的 Gröbner 基。若采用临界对删减准则，也应给出相应判据或补回被省去的验证。

\ProseParagraphBreak
保存表示矩阵并不要求每次重新求解理想成员问题。初始输入附单位向量；每次乘单项式、相减和归一化时，对表示向量作同样运算，即可始终保存 \(g=Fu\)。F4 的行操作可以同样作用于表示矩阵。FGLM 返回的每个关系可先用旧 Gröbner 基求余并取得表示，再与旧 Gröbner 基到输入的变换矩阵相乘。

\begin{example}\label{b3:exm:B-complete-certificate}
只检查 \(F\) 对 \(G\) 的余项为零是不够的：任意输入对 \(G=\{1\}\) 都有零余项，但输入理想未必为整个环。反过来，只检查 \(G\subset(F)\) 也不够：空集或全零列不会恢复非零输入理想。两个方向的矩阵等式正是排除这两种错误的必要信息。

下面给出一份包含两向表示和临界对记录的完整证书。设 \(k\) 为域，\(S=k[x,y]\) 采用 \(x>y\) 的字典序，取
\[
F=(f_1,f_2)=(x^2-y,xy-1),\qquad
G=(g_1,g_2)=(x-y^2,y^3-1).
\]
给出的表示矩阵为
\[
U=\begin{pmatrix}y&-y^2\\-x&1+xy\end{pmatrix},\qquad
V=\begin{pmatrix}x+y^2&y\\y&1\end{pmatrix}.
\]
两条矩阵等式可逐列展开核验：
\[
\begin{aligned}
yf_1-xf_2&=(x^2y-y^2)-(x^2y-x)=g_1,\\
-y^2f_1+(1+xy)f_2&=(-x^2y^2+y^3)+(x^2y^2-1)=g_2,\\
(x+y^2)g_1+yg_2&=(x^2-y^4)+(y^4-y)=f_1,\\
yg_1+g_2&=(xy-y^3)+(y^3-1)=f_2.
\end{aligned}
\]
因此 \(G=FU\)、\(F=GV\)，两组生成元给出同一理想。两个首单项式分别为 \(x,y^3\)，只有一个临界对，其最小公倍数为 \(xy^3\)，且
\[
S(g_1,g_2)=y^3g_1-xg_2=x-y^5=g_1-y^2g_2.
\]
从 \(x-y^5\) 减去 \(g_1\)，余下 \(-y^5+y^2=-y^2g_2\)，再约化为零。标准表示右端的两个乘积具有首单项式 \(x\) 和 \(y^5\)，而在此字典序下 \(y^5<x<xy^3\)，所需严格上界成立。由\cref{b3:thm:buchberger-criterion}，\(G\) 是输入理想的 Gröbner 基。整个核验只用两条矩阵等式、首项和这一份标准表示，无须重新运行生成 Gröbner 基的算法。
\end{example}

成员证明则更简单：只要明确写出 \(f=\sum a_if_i\)，即可逐项展开核验，不需要先证明给定生成组为 Gröbner 基。非成员结论需要更多结构，例如对已经验证的 Gröbner 基取得非零正规形。答案的生成方法可以很复杂，而最后检查这些等式的过程可以保持直接。

% !TeX root = ../../Book3.tex
% 纲要标题：### B.3 分解、参数与组合

\section{分解、参数与组合}
\label{b3:sec:B03}

一个分解或正规化的候选输出，必须同时满足包含关系、终止条件与所声称的结构性质。参数化方程还要求说明输出在哪些参数取值下成立，不能把一次泛型计算当作全部纤维的答案。

% 纲要细目（与计划纲要同步）：
% * 准素分解、根理想、饱和与正规化的有效问题
% * 参数化方程的综合 Gröbner 系统、正则链及参数空间分层
% * 一般参数值的答案不代表所有特殊化；退化参数需要独立分支
% * 单项式理想、组合结构与 Hilbert 数据；分解结果的唯一性条件
%
% #### B.3.1 准素分解、根理想、饱和与正规化的有效问题
% #### B.3.2 综合 Gröbner 系统、正则链与参数空间分层
% #### B.3.3 一般参数、特殊化与退化分支
% #### B.3.4 单项式理想、组合结构与分解的非唯一性


\subsection{准素分解、根理想、饱和与正规化的有效问题}
\label{b3:subsec:B03:01}

正文已经证明，饱和、根理想、准素分解和整闭包具有各自的代数意义。有效计算的任务是从有限输入得到有限输出，并附上足以识别该输出的依据。四种输出不能混为一谈：根理想舍去幂零结构；准素分解保存理想本身；饱和选择某些分支；正规化则通常改变坐标环而非只把原理想换成另一个根理想。

\begin{example}\label{b3:exm:B-decomposition}
在 \(k[x,y]\) 中，
\[
I=(x^2,xy)=(x)\cap(x^2,y),\qquad
\sqrt I=(x),\qquad I:y^\infty=(x).
\]
交式可逐项检查：若 \(xh\in(x^2,y)\)，模 \(y\) 后 \(h(x,0)\) 被 \(x\) 整除，所以 \(h\in(x,y)\)。第二个理想的商为 \(k[x]/(x^2)\)，故它是 \((x,y)\)-准素的。饱和的等式来自 \(yx\in I\) 及 \(I\subset(x)\)，其中 \(y\notin(x)\)。这个例子中的饱和移除了嵌入的加厚结构，而根理想同时忘记了全部幂零信息。
\end{example}

对于域上的有限多项式输入，交和饱和可用\cref{b3:prop:ideal-intersection-elimination,b3:prop:colon-saturation-elimination}的辅助变量构造化为消去。核验一个候选准素分解时，还要证明各分量准素、交等于原理想，以及在声称不可约简时没有冗余。仅列出若干零点集的并不够，因为它不能区分 \(I\) 与 \(\sqrt I\)。

\ProseParagraphBreak
正规化的输出则应包含有限代数及其到全分式环的嵌入。例如 \(A=k[t^2,t^3]\subset k[t]\) 的正规化为 \(k[t]\)：\(t\) 满足首一方程 \(T^2-t^2=0\)，且 \(k[t]=A+At\)；两者分式域相同，而 \(k[t]\) 是整闭环。若一个分式对 \(A\) 整，则同一首一方程也说明它对 \(k[t]\) 整，因此它属于 \(k[t]\)。这核验了该例；一般正规化算法还须处理有限性假设、分支及系数域。

\subsection{综合 Gröbner 系统、正则链与参数空间分层}
\label{b3:subsec:B03:02}

参数计算的难处不是增加几个字母，而是某些系数随参数取值变为零。把参数当作有理函数域的元素，可得到一般情形的答案；要覆盖全部参数，还须记录除法中要求哪些系数非零。

\begin{definition}\label{b3:def:B-comprehensive-system}
设 \(k\) 为域，\(\overline k\) 为其代数闭包，\(a=(a_1,\ldots,a_r)\) 为参数，\(x=(x_1,\ldots,x_n)\) 为未知量，\(F\subset k[a,x]\) 为有限集，并固定 \(x\) 的全局项序。若有限个由参数多项式等式和非等式定义的集合 \(C_\nu\subset\overline{k}^{\,r}\) 覆盖参数空间，并各附一列可在 \(C_\nu\) 上有定义的多项式 \(G_\nu\)，使对每个 \(c\in C_\nu\)，\(G_\nu(c,x)\) 是 \((F(c,x))\) 的 Gröbner 基，则这些分支数据称为一个\term{综合 Gröbner 系统}[comprehensive Gröbner system]。
\end{definition}

这一定义描述输出应满足什么，不声称任意一次一般参数计算已经给出这样的覆盖。若分支重叠，其特殊化答案须描述同一输入理想；若要求互不相交，还需进一步细分条件。

\ProseParagraphBreak
三角形方程组提供另一种组织方式。按变量 \(x_1<\cdots<x_n\) 排列，一个非常数多项式实际出现的最高变量称为它的主变量；把它视为主变量的多项式时，其首系数称为初式。

\begin{definition}\label{b3:def:B-regular-chain}
设 \(K\) 为域，\(T=(t_1,\ldots,t_s)\subset K[x_1,\ldots,x_n]\) 为非常数多项式列，主变量严格递增，初式依次为 \(h_1,\ldots,h_s\)。令 \(H_{i-1}=h_1\cdots h_{i-1}\)，并约定空乘积为 \(1\)。若每个 \(h_i\) 在
\[
K[x_1,\ldots,x_n]/
\bigl((t_1,\ldots,t_{i-1}):H_{i-1}^{\infty}\bigr)
\]
中均为非零因子，且这些商环非零，则称 \(T\) 为一条\term{正则链}[regular chain]。
\end{definition}

这个条件容许前面的商环有零因子，却排除所需初式成为零因子。例如 \((x^2,xy+1)\) 关于 \(x<y\) 是三角形方程组，但第二个初式 \(x\) 在 \(K[x,y]/(x^2)\) 中为零因子，故不是正则链。链所描述的基本集合通常是
\[
V(T)\setminus V(h),\qquad
h=\text{各初式的乘积},
\]
其闭包与饱和理想 \((T):h^\infty\) 有关。因而“正则链”既不等于任意三角生成组，也不自动等于素理想的一组生成元。一般正则链分解的构造和终止证明属于参数计算专门理论，这里用下面的分支计算说明为何需要它们。

\subsection{一般参数、特殊化与退化分支}
\label{b3:subsec:B03:03}

\begin{example}\label{b3:exm:B-parameter-branches}
取参数 \(a,b\)，未知量 \(x,y\)，方程组
\[
ax-b=0,\qquad ay=0.
\]
在参数域 \(k(a,b)\) 上得到 \(x=b/a,\ y=0\)。但完整参数空间必须分为
\[
\begin{array}{c|c|c}
\text{参数条件}&\text{特殊化理想}&\text{一个 Gröbner 基}\\ \hline
a\ne0&(x-b/a,y)&\{x-b/a,y\}\\
a=0,\ b\ne0&(1)&\{1\}\\
a=0,\ b=0&(0)&\varnothing
\end{array}
\]
三种情况分别有一个点、无解以及整个仿射平面。各行均由代入和可逆系数的除法直接核验，并且这些条件互不相交、覆盖全部参数。
\end{example}

一般参数答案没有“错误”；它准确描述开集 \(a\ne0\)。错误在于省去这个条件，把出现的分母形式上约掉以后当作全体参数的结论。\cref{b3:prop:groebner-specialization-open}给出更一般的安全情形：有限计算所用的系数条件同时不退化时，证书可以特殊化。

\ProseParagraphBreak
有些退化不改变解集的点数，却改变重数。例如 \(x^2-a\) 在代数闭域上、特征不为 \(2\) 时，\(a\ne0\) 有两个不同根，\(a=0\) 只有一个加厚的根。两个商代数的向量空间维数都为 \(2\)，所以“维数保持”不能代替根理想或可约性判定。把参数分支与商代数的乘法矩阵一起记录，可以看出简单特征值怎样合并为带幂零部分的局部代数。

\subsection{单项式理想、组合结构与分解的非唯一性}
\label{b3:subsec:B03:04}

单项式理想把成员关系化为指数向量的比较。设
\[
J=(x^{\alpha^{(1)}},\ldots,x^{\alpha^{(s)}})\subset
S=k[x_1,\ldots,x_n].
\]
单项式 \(x^\beta\) 属于 \(J\)，当且仅当某个 \(\alpha^{(i)}\) 在每个坐标上都不超过 \(\beta\)。标准单项式的计数因而可以用有限个指数区域的并来组织。

\begin{remark}\label{b3:prop:B-monomial-hilbert}
单项式理想的 Hilbert 级数计算直接使用\cref{b3:prop:hilbert-initial-count}已经证明的容斥公式：对生成元的每个子集计算最小公倍式的次数，再按子集大小交错相加，所得分子除以 \((1-t)^n\)。实际所需数据只有这些最小公倍式及其次数，不必先计算自由消解。多个子集若有相同的最小公倍式，可以合并其带符号的贡献；冗余生成元虽不影响答案，却会增加需要枚举的子集数。
\end{remark}

例如取 \(J=(x^3,x^2y,y^3)\subset k[x,y]\)。三个生成元的次数都为 \(3\)，两两最小公倍式的次数依次为 \(4,6,5\)，三个生成元的最小公倍式次数为 \(6\)。所以两个 \(t^6\) 项抵消，得到
\[
H_{k[x,y]/J}(t)
=\frac{1-3t^3+t^4+t^5}{(1-t)^2}
=1+2t+3t^2+t^3.
\]
也可以按 \(x\) 的指数分组：指数零和一时，各允许 \(y^0,y,y^2\)；指数二时只允许 \(y^0\)。由此得到七个标准单项式
\[
1,\ y,\ y^2,\ x,\ xy,\ xy^2,\ x^2,
\]
逐次计数与上面的级数一致，并给出商的长度为 \(7\)。这个例子同时核验了最小公倍式的容斥计算与直接的指数计数。

\ProseParagraphBreak
准素分解还有另一种非唯一性。前例 \((x^2,xy)\) 对每个整数 \(r\ge1\) 都有
\[
(x^2,xy)=(x)\cap(x^2,xy,y^r).
\]
第二个商环中 \(x,y\) 都幂零，且剩余域为 \(k\)，故第二分量为 \((x,y)\)-准素；交式由单项式整除即可核验。这些嵌入准素分量随 \(r\) 改变。关联素理想集合具有正文已经说明的唯一性，而具体嵌入分量并无同样唯一性。比较两个程序输出时，应比较它们是否满足所要求的分解性质，不能只因列表不同就判定其中一个错误。

% !TeX root = ../../Book3.tex
% 纲要标题：### B.4 模与自由消解的计算
% 纲要位置：工作记录/计划纲要.md:2853

\section{模与自由消解的计算}
\label{b3:sec:B04}

关系矩阵随生成元的更换而改变，计算记录必须保存这些更换的双向表示。尤其在消解中，相邻矩阵乘积为零只给出像包含于核，核的生成证书才证明正合性。

% 纲要内容（仅作写作计划，尚非教材正文）：
% * 调用第21章的 Schreyer 构造与分次极小化，以证书对照表区分全部核、复形、正合性与 Betti 数
% * 局部标准基方法归附录 F；本节证明局部成员与原环中乘性集倍数证书的等价性
% * SINGULAR 等精确代数系统可作实验载体；工作环、项序与局部化信息必须记录
% * 多项式与局部环中的同一输入可能有不同的理想成员答案，应提供对照实例
% #### B.4.1 Schreyer 构造、自由消解与 Betti 表
% #### B.4.2 局部计算中的项序与结果验证
% #### B.4.3 计算机代数系统中的环与局部化
% #### B.4.4 多项式环与局部环中理想成员判定的比较


\subsection{Schreyer 构造、自由消解与 Betti 表}
\label{b3:subsec:B04:01}

设 \(k\) 为域，\(S=k[x_1,\ldots,x_n]\)，输入矩阵 \(F\in S^{r\times s}\) 的列生成子模 \(N\subset S^r\)。计算模 Gröbner 基矩阵 \(G\in S^{r\times t}\) 时，保存双向表示
\[
G=FU,\qquad F=GV,
\quad U\in S^{s\times t},\quad V\in S^{t\times s}.
\]
由\cref{b3:thm:schreyer-syzygies}求出生成 \(\ker G\) 的矩阵 \(T\) 后，\cref{b3:alg:polynomial-matrix-kernel}将原核还原为
\[
\ker F=\operatorname{im}K,\qquad
K=\bigl[\,UT\ \ I_s-UV\,\bigr].
\]
第二块记录原生成组与新生成组之间多余的表示；它可能含有第一块没有给出的关系。

\begin{example}\label{b3:exm:B-generator-change}
设 \(k\) 为域，在 \(S=k[x,y]\) 中取 \(F=(x,y,x+y)\)，改用 \(G=(x,y)\)。可取
\[
U=\begin{pmatrix}1&0\\0&1\\0&0\end{pmatrix},
\quad
V=\begin{pmatrix}1&0&1\\0&1&1\end{pmatrix},
\quad T=\binom{y}{-x}.
\]
因此原关系由
\[
(y,-x,0)^{\mathsf T},\qquad(-1,-1,1)^{\mathsf T}
\]
生成。第二个关系来自被删掉的冗余生成元，单算新基的合冲而不记录变换，就会漏掉它。
\end{example}

同样的记录要逐层保存。\cref{b3:tab:B-resolution-certificates}列出各项结论所需的数据，其中 \(C_\bullet\xrightarrow{\varepsilon}M\) 是有限秩自由模的增广序列，\(d_i:C_i\rightarrow C_{i-1}\)；讨论分次极小性与 Betti 数时，\(S\) 取标准分次，\(M\) 有限生成，各映射齐次且次数为零。

\begin{table}[htbp]
\centering
\caption{合冲、消解与 Betti 数的计算证书}
\label{b3:tab:B-resolution-certificates}
\begin{tabularx}{\linewidth}{@{}>{\raggedright\arraybackslash}p{0.20\linewidth}>{\raggedright\arraybackslash}X@{}}
\toprule
所需结论 & 必须保存或核验的数据 \\
\midrule
生成子模相等 & \(G=FU\) 与 \(F=GV\) 的精确矩阵等式 \\
全部核已求出 & 已验证的模 Gröbner 基 \(G\)、Schreyer 合冲矩阵 \(T\) 及两向变换 \(U,V\)；由还原公式得到全部核生成矩阵 \(K\)，不只检查 \(FK=0\) \\
序列为复形 & 各层 \(d_i d_{i+1}=0\)，且 \(\varepsilon d_1=0\) \\
复形为消解 & \(\varepsilon\) 满射、\(\ker\varepsilon=\operatorname{im}d_1\)，以及每一层的全部核与下一层像相等；有限左端还须核验最左侧微分单射 \\
分次消解极小 & 已验证正合的消解，其微分项全在 \(\mathfrak m=(x_1,\ldots,x_n)\) 中；消去单位项时须同时更新相邻矩阵 \\
Betti 数 & 记录极小消解中 \(S(-j)\) 在 \(C_i\) 中的重数 \(\beta_{i,j}\)，它等于 \(\dim_k\operatorname{Tor}^S_i(M,k)_j\) \\
\bottomrule
\end{tabularx}
\end{table}

\cref{b3:thm:hilbert-syzygy}给出有限生成分次模具有有限长度自由消解的依据，\cref{b3:prop:graded-cancellation}给出单位项消去的具体步骤；极小消解与 Tor 的关系见\eqref{b3:eq:graded-betti-tor}。这些结论使求核与极小化各有明确的核验目标。

\subsection{局部计算中的项序与结果验证}
\label{b3:subsec:B04:02}

全局模项序的首项约化以良序为终止依据。局部序允许 \(1>x>x^2>\cdots\)，会改变可逆元素，也会改变约化方法。因此从多项式环改为局部计算时，不能仅把排序参数换掉，再沿用原来的除法终止证明。

\ProseParagraphBreak
\appref{b3:app:F}将定义局部标准基和弱正规形。局部成员关系可以完全写成原环中的有限等式。

\begin{proposition}[局部成员证书]\label{b3:prop:B-local-membership-certificate}
设 \(S\) 为交换环，\(U\subset S\) 为不含零的乘性集，\(R=U^{-1}S\)。给定 \(F\in S^{r\times s}\) 与 \(f\in S^r\)，则
\[
\dfrac{f}{1}\in\operatorname{im}(U^{-1}F:R^s\longrightarrow R^r)
\quad\Longleftrightarrow\quad
\text{存在 }u\in U,\ a\in S^s\text{，使 }uf=Fa.
\]
\end{proposition}
\begin{proof}
若 \(uf=Fa\)，在 \(R\) 中除以 \(u\) 即得成员表示。反过来，若 \(f/1\) 属于像，将表示中有限多个分母通分，可写成 \(f/1=Fb/v\)，其中 \(v\in U\)、\(b\in S^s\)。局部化中的相等意味着存在 \(w\in U\) 使 \(w(vf-Fb)=0\)。取 \(u=wv\)、\(a=wb\)，便有 \(uf=Fa\)。此处乘上 \(w\) 的步骤也涵盖 \(S\) 有零因子的情形。
\end{proof}

因此验证成员证书时，先检查原环中的矩阵等式，再核验 \(u\) 属于指定乘性集。在 \(S=k[x_1,\ldots,x_n]\)、\(U=S\setminus(x_1,\ldots,x_n)\) 的原点局部化中，后一个条件就是 \(u(0)\ne0\)。

\ProseParagraphBreak
非成员证明需要一个已验证的局部标准基及适用的正规形定理。仅有一次未约化为零的尝试不能作数，因为约化方法可能不完整。若输出为局部自由消解，每层仍要说明全部核的生成，并在作单位消去后再次检查相邻微分。局部极小性由极大理想中的微分项判断，而不是由多项式总次数或文件中列数的变化判断。

\subsection{计算机代数系统中的环与局部化}
\label{b3:subsec:B04:03}

一份可复算的环描述至少包括系数域、变量、变量次序、项序、商关系以及被倒置的元素。若系数是代数扩张，还要写出其定义多项式；若对象带分次或滤过，还要记录各变量和基向量的次数。只保存一张系数表无法恢复这些信息。

\ProseParagraphBreak
例如以下三个工作环有同名变量，却表达不同的问题：
\[
\mathbb Q[x,y],\qquad
\mathbb Q[x,y]_{(x,y)},\qquad
\mathbb Q[\![x,y]\!].
\]
第一环中的多项式 \(1-x\) 不可逆，后两个环中则可逆。第二环的元素是满足分母条件的有理函数；第三环还允许一般形式级数，所以它们并非相同的环。使用有限截断计算第三环时，还必须附上精度及为何该精度足够的说明。

\ProseParagraphBreak
SINGULAR 等精确计算系统可以承载这些实验，但报告不应以某个命令名称代替数学定义。应把“输入理想、输出标准基、表示矩阵、余项或 Hilbert 数据”分别保存。软件内部使用何种临界对策略可以更换；公开的结果仍应由明确的环、项序和有限等式决定。\subsecref{b3:subsec:F04:01}将给出截断何时成为完整长度证书的具体条件。

\subsection{多项式环与局部环中理想成员判定的比较}
\label{b3:subsec:B04:04}

\begin{example}\label{b3:exm:B-local-membership}
设 \(k\) 为域，\(S=k[x]\)，\(I=(x-x^2)\)。在 \(S\) 中，\(x\notin I\)：若 \(x=(x-x^2)h\)，整环中约去 \(x\) 得 \(1=(1-x)h\)，与 \(1-x\) 非单位矛盾。但在原点局部环 \(S_{(x)}\) 中，
\[
x=(1-x)^{-1}(x-x^2),
\]
故 \(IS_{(x)}=(x)\)。在点 \(x=1\) 的局部环中，反而是 \(x\) 可逆，于是 \(IS_{(x-1)}=(1-x)\)。
\end{example}

从商环看，这个差别同样清楚。多项式商 \(k[x]/(x-x^2)\) 有两个点，并由中国剩余定理同构于 \(k\times k\)；在原点局部化只保留一个分量。局部成员关系的改变对应实际删去远离所选点的部分，不是排序技巧造成的近似。

\ProseParagraphBreak
因此，计算前首先应决定研究的是整个代数集、指定点附近的局部环，还是完备局部环。全局求解可以先计算零维商代数，再按局部分量分析；局部标准基也可以直接计算指定点的结构。两种方法能够在适当比较映射下核对，但它们首先必须确实描述同一个局部对象。


\begin{chaptersummary}
有限输入首先要说明它表示哪一个环、理想或模。系数能否精确判零、首系数能否求逆、分母是否属于指定乘性集，都会改变算法的合法步骤。FGLM 的有限性来自商代数的有限维数：候选单项式中线性无关的类至多出现 \(D\) 个，其余类给出新项序的关系。

\ProseParagraphBreak
计算方法与验证方法可以分开。行消元、签名、模素数重构和参数分支有助于产生答案；两向生成元表示、首项判据和精确矩阵等式使答案能够独立复核。模的计算还必须保存生成元变换，并在每一层证明像等于核。单项式计数与 Hilbert 数据很有用，但不能代替理想相等、准素性或消解正合性的证明。
\end{chaptersummary}

\BookExercises
以下各题的系数运算均为精确运算；未另作说明时，\(k\) 为域。

\begin{exercise}\label{b3:ex:B-truncation}
设 \(k\) 为域，\(N\ge1\)。只给定 \(f\in k[\![x]\!]\) 模 \((x)^N\) 的截断，能否判断 \(f=0\)，或在 \(f\ne0\) 时确定其阶？若截断中的最低非零项为 \(a_rx^r\)，其中 \(r<N\)，可以确定什么？
\end{exercise}
\begin{solution}
零截断同时允许 \(f=0\)、\(f=x^N\) 和 \(f=x^{N+1}\)，故不能判断是否为零，也不能确定非零时的阶。若截断已出现最低非零项 \(a_rx^r\)，后面的未知项次数均至少为 \(N>r\)，不能改变该项，故可确定 \(f\ne0\)、\(\operatorname{ord}(f)=r\) 以及初始形式 \(a_rx^r\)。
\end{solution}

\begin{exercise}\label{b3:ex:B-strong-basis}
在 \(\mathbb Z[x]\) 中证明 \(\{2,x\}\) 是理想 \(I=(2,x)\) 的强 Gröbner 基。求 \(\mathbb Z[x]/I\)，并说明扩张到 \(\mathbb Q[x]\) 后发生什么。
\end{exercise}
\begin{solution}
属于 \(I\) 的常数为偶数，所以次数零的非零首项被 \(2\) 整除；次数至少为一的首项都被 \(x\) 整除。故给定集合是强 Gröbner 基。商环先令 \(x=0\)，再模 \(2\)，得到 \(\mathbb F_2\)。在 \(\mathbb Q[x]\) 中 \(2\) 可逆，扩张理想为整个环，商为零环；因此从分式域上的零商不能恢复原商的整数扭结构。
\end{solution}

\begin{exercise}\label{b3:ex:B-fglm-nonreduced}
设 \(k\) 为域。对 \(I=(y-x^2,x^3)\subset k[x,y]\)，已知 \(y>x\) 字典序的标准单项式基为 \(1,x,x^2\)。用 FGLM 改为 \(x>y\) 字典序，写出候选的处理次序、返回基及标准单项式。解释为什么算法不需要 \(I\) 为根理想。
\end{exercise}
\begin{solution}
在旧商环中 \(y=x^2\)、\(x^3=0\)。先接受 \(1\)，候选为 \(y,x\)；接受 \(y\)，其坐标是 \(x^2\)。接着处理 \(y^2\)，坐标为零，产生 \(y^2\)。随后接受 \(x\)，坐标与 \(1,x^2\) 独立；处理 \(xy\) 时坐标为 \(x^3=0\)，产生 \(xy\)；处理 \(x^2\) 时坐标等于 \(y\)，产生 \(x^2-y\)。返回
\[
G=(y^2,xy,x^2-y),\qquad B=(1,y,x).
\]
这些首项为 \(y^2,xy,x^2\)，商维数为三。\(x\) 在商中是非零幂零元，所以 \(I\) 非根理想；FGLM 使用的是坐标线性相关和有限维数，从未要求乘法矩阵可对角化。
\end{solution}

\begin{exercise}\label{b3:ex:B-fglm-unit}
设 \(k\) 为域，\(a,b\ge1\)，对 \(I=(x^a,y^b)\subset k[x,y]\) 使用 FGLM 的只从接受单项式产生后继的队列，并取 \(x>y\) 的字典序。求接受的标准单项式集合与被处理过的全部不同候选数，将精确计数同 \(1+2\dim_k(k[x,y]/I)\) 比较。
\end{exercise}
\begin{solution}
接受集合为 \(B=\{x^iy^j:0\le i<a,\ 0\le j<b\}\)，有 \(D=ab\) 个元素。所有候选都是 \(\{1\}\cup xB\cup yB\)：除这块矩形内的单项式外，还有 \(x^ay^j\)（\(0\le j<b\)）和 \(x^iy^b\)（\(0\le i<a\)）两条边上的候选，故精确个数为 \(ab+a+b\)。边上的候选给出关系或被已有首项丢弃，均不再扩展。由 \((a-1)(b-1)\ge0\)，有 \(ab+a+b\le1+2ab\)。一般界可以把不同接受项产生的同一候选重复计数，这里去重后的精确值通常更小。
\end{solution}

\begin{exercise}\label{b3:ex:B-rational-reconstruction}
模 \(101\) 的剩余类 \(34\) 在范围 \(|a|\le5,\ 1\le b\le5\) 内重构为互素分数 \(a/b\)，并要求 \(b\) 模 \(101\) 可逆。求答案并证明唯一。
\end{exercise}
\begin{solution}
\(3\cdot34=102\equiv1\pmod{101}\)，故 \(a/b=1/3\) 满足条件。若有另一解 \(c/d\)，交叉差 \(ad-bc\) 被 \(101\) 整除，而其绝对值至多 \(2\cdot5\cdot5=50<101\)，故它等于零。互素且分母为正的分数表示唯一。
\end{solution}

\begin{exercise}\label{b3:ex:B-one-sided-certificate}
给出候选生成组 \(G\subset(F)\) 但 \((G)\ne(F)\) 的例子，以及 \((F)\subset(G)\) 但二者不等的例子。说明哪些附加数据可以排除它们。
\end{exercise}
\begin{solution}
在 \(k[x]\) 中取 \(F=(x)\)、\(G=(x^2)\)，有 \(G=F(x)\)，但 \(x\notin(x^2)\)。反向取 \(F=(x^2)\)、\(G=(x)\)，则 \(F=G(x)\)，但 \(x\notin(F)\)。同时给出 \(G=FU\) 与 \(F=GV\) 的矩阵等式便排除这两类遗漏；要进一步声称 Gröbner 基，还需验证临界对或等价的首项条件。
\end{solution}

\begin{exercise}\label{b3:ex:B-complete-certificate}
设 \(k\) 为域，\(S=k[x,y]\) 采用 \(x>y\) 的字典序，\(F=(x^2-y,xy)\)。构造首一 Gröbner 基 \(G\)，给出两向表示矩阵及全部临界对的标准表示。再据此求标准单项式基，证明商环同构于 \(k[x]/(x^3)\)，从而区别唯一底层点与商空间的三维长度。
\end{exercise}
\begin{solution}
可取 \(G=(g_1,g_2,g_3)=(x^2-y,xy,y^2)\)，以及
\[
U=\begin{pmatrix}1&0&-y\\0&1&x\end{pmatrix},\qquad
V=\begin{pmatrix}1&0\\0&1\\0&0\end{pmatrix}.
\]
直接展开可核验 \(G=FU\)、\(F=GV\)。三个临界对为
\[
S(g_1,g_2)=-g_3,\qquad
S(g_1,g_3)=-yg_3,\qquad S(g_2,g_3)=0.
\]
前两式右端的首单项式为 \(y^2,y^3\)，分别严格小于 \(x^2y,x^2y^2\)，所以确为标准表示，Buchberger 判据适用。首项理想为 \((x^2,xy,y^2)\)，标准单项式为 \(1,x,y\)。关系 \(y=x^2\)、\(xy=0\) 又给出 \(S/(F)\cong k[x]/(x^3)\)，其中 \(x\ne0\) 且 \(x^3=0\)。故只有一个底层点，商的 \(k\)-维数与合成长度却均为 \(3\)。
\end{solution}

\begin{exercise}\label{b3:ex:B-parameter}
设 \(k\) 为域，\(a,b\) 为参数。对 \(a x^2+b x\in k[a,b,x]\)，给出覆盖全部参数的分支及各分支的首一 Gröbner 基。在 \(a\ne0\) 的分支内，再比较 \(b=0\) 与 \(b\ne0\) 两种情形的商代数维数与根结构。
\end{exercise}
\begin{solution}
分为 \(a\ne0\)、\(a=0,b\ne0\)、\(a=b=0\)。相应基为
\[
\{x^2+(b/a)x\},\qquad \{x\},\qquad\varnothing.
\]
商维数依次为 \(2,1,\infty\)。在第一分支内，\(b\ne0\) 时有两个不同根 \(0,-b/a\)，\(b=0\) 时只有一个二重根；商维数同为二。此结论在任意特征成立，因为 \(b/a\ne0\) 就使两个根不同。
\end{solution}

\begin{exercise}\label{b3:ex:B-primary-components}
设 \(k\) 为域。证明对每个 \(r\ge1\)，
\[
(x^2,xy)=(x)\cap(x^2,xy,y^r)
\]
是 \(k[x,y]\) 的一个不可删减准素分解。
\end{exercise}
\begin{solution}
交理想的单项式生成元由两组生成元的最小公倍数给出，即 \(x^2,xy,xy^r\)，最后一个被 \(xy\) 整除，得到原理想。\((x)\) 为素理想。第二分量的商中 \(x,y\) 幂零，且其唯一极大理想为 \((x,y)\) 的像；极大理想外的元素为单位，故零因子只能是幂零元，第二分量准素。\((x)\) 不包含于第二分量，因为 \(x\) 不属于它；第二分量也不包含于 \((x)\)，因为 \(y^r\notin(x)\)。因此任何分量都不可删去。
\end{solution}

\begin{exercise}\label{b3:ex:B-monomial-series}
设 \(k\) 为域。求 \(k[x,y]/(x^3,x^2y,y^2)\) 的标准单项式、Hilbert 级数和长度，并用容斥公式核对。
\end{exercise}
\begin{solution}
标准单项式为 \(1,x,x^2,y,xy\)，故级数为 \(1+2t+2t^2\)，长度为 \(5\)。容斥分子为
\[
1-(t^3+t^3+t^2)+(t^4+t^5+t^4)-t^5
=1-t^2-2t^3+2t^4.
\]
它等于 \((1-t)^2(1+2t+2t^2)\)，与标准单项式计数一致。
\end{solution}

\begin{exercise}\label{b3:ex:B-module-change}
设 \(k\) 为域，\(S=k[x,y]\)，\(F=(x,y,x+y):S^3\rightarrow S\)。直接证明其全部关系由 \((y,-x,0)^{\mathsf T}\) 与 \((-1,-1,1)^{\mathsf T}\) 生成，不使用 Gröbner 基。
\end{exercise}
\begin{solution}
若 \(ax+by+c(x+y)=0\)，则 \((a+c)x+(b+c)y=0\)。因 \(x,y\) 互素，存在 \(h\) 使 \(a+c=hy,\ b+c=-hx\)。所以
\[
(a,b,c)^{\mathsf T}
=h(y,-x,0)^{\mathsf T}+c(-1,-1,1)^{\mathsf T}.
\]
两个给定向量显然都是关系，故生成全部核。
\end{solution}

\begin{exercise}\label{b3:ex:B-local-vs-global}
设 \(k\) 为域，在 \(k[x]\) 中取 \(I=(x(x-1)^2)\)。分别求其在 \((x)\) 与 \((x-1)\) 处局部化后的理想，以及两个局部商的长度。
\end{exercise}
\begin{solution}
在 \((x)\) 处，\(x-1\) 可逆，故扩张理想为 \((x)\)，商为 \(k\)，长度一。在 \((x-1)\) 处，\(x\) 可逆，故扩张理想为 \(((x-1)^2)\)，商基为 \(1,x-1\)，长度二。全局商有维数三，与两个局部长度之和一致。
\end{solution}

\begin{exercise}\label{b3:ex:B-complex-not-resolution}
设 \(k\) 为域，\(S=k[x,y]\)。考虑
\[
0\longrightarrow S\xrightarrow{\,d_2\,}S^2
\xrightarrow{\,d_1\,}S\longrightarrow S/(x,y)\longrightarrow0,
\qquad
d_2=\begin{pmatrix}-xy\\x^2\end{pmatrix},\quad d_1=\begin{pmatrix}x&y\end{pmatrix}.
\]
验证复形条件，求 \(S^2\) 处的同调模，并修改 \(d_2\) 得到真正的自由消解。
\end{exercise}
\begin{solution}
直接相乘得 \(d_1d_2=-x^2y+yx^2=0\)。若 \(xu+yv=0\)，因 \(x,y\) 互素，可写 \(u=-yh\)、\(v=xh\)，所以 \(\ker d_1=S(-y,x)^{\mathsf T}\)。而 \(\operatorname{im}d_2=xS(-y,x)^{\mathsf T}\)，该处上同调为 \(S/(x)\ne0\)。左端映射因 \(S\) 是整环而单射，右端也正合；唯一缺口正是在中间。把 \(d_2\) 改成列 \((-y,x)^{\mathsf T}\)，其像便等于已求出的核，且仍单射，从而得到完整的自由消解。
\end{solution}
